\documentclass [prd,reprint,eqsecnum,nofootinbib]  {revtex4-1}

\usepackage{amsfonts,amsmath,amscd,mathrsfs,hyperref,xcolor,ulem}
\normalem
\newcommand{\corr}[1]{\textcolor{blue}{#1}}
\newcommand{\add}[1]{\textcolor{red}{#1}}
\newcommand{\del}[1]{\textcolor{green}{\sout{#1}}}

\def\SO{\mathsf{SO}}

\def\so{\mathsf{so}}

\def\SU{\mathsf{SU}}

\def\cP{{\cal P}}

\def\cD{{\cal D}}

\def\cC{{\cal C}}

\def\cL{{\cal L}}

\def\w{\omega}

\def\v{\wedge}

%\usepackage[colorinlistoftodos, shadow]{todonotes} %notatki: paczka

%Komendy do stosowania:

%\todo{Comment on margin with line}

%\todo[inline]{Comment in text}

%\todo[noline]{Comment on margin without line}

%\missingfigure{Missing picture here.}

%\listoftodos %notatki: lista rzeczy do poprawy



\begin{document}

\title{%

Local Maxwell symmetry and gravity}

\author{Professor Gerald Lambeau}
\affiliation{Department of Mathematics, Massachusetts Institute of Technology, Cambridge, Massachusetts, USA}
\author{Will Hunting}
\affiliation{Independent Scholar, Boston, Massachusetts, USA}
\date{\today}

\begin{abstract}\noindent
\corr{We examine several physical aspects of a theory obtained by
gauging the AdS--Maxwell symmetry. The resulting theory takes the form
of Einstein gravity coupled to an $\sf{SO(3,1)}$ Yang--Mills field.
We identify a second tetrad field associated with a linear combination
of the Lorentz and Maxwell connections. Treating this tetrad as a
fundamental variable allows us to recast the theory as a bi-metric
($f$--$g$) theory of gravity in the first-order formalism. Finally,
we study a simple cosmological model derived from AdS--Maxwell gravity.}
\end{abstract}



\maketitle
\add{The discussion is deliberately limited to the classical field-theoretic
and homogeneous cosmological settings specified below. In particular, the
formal similarities with familiar gauge and bi-metric constructions should
not be taken, by themselves, as evidence for quantum consistency or
phenomenological viability. Those questions require separate analyses.}



\section{Introduction}

\corr{The Maxwell symmetry extends Poincar\'e symmetry and describes
fields in a constant electromagnetic background}
\cite{Bacry:1970ye}, \cite{Schrader:1972zd}. This symmetry is of
interest \corr{from a purely algebraic point of view} because it
circumvents \corr{a well-known theorem} that precludes a central
extension of the Poincar\'e algebra (see, e.g.,
\cite{Galindo:1967}, \cite{Soroka:2011tc}, \cite{Gomis:2009dm}).

\corr{Until recently, however, the Maxwell symmetry attracted
relatively little attention. This is somewhat surprising because
physical systems in constant electromagnetic backgrounds, or in
closely related settings, are frequently encountered and play an
important role in physics. This suggests that Maxwell symmetry may
have an important role to play in physical theories. For example,
since canonical noncommutativity (for a review, see
\cite{Douglas:2001ba}) is ultimately related to the behaviour of
quantum systems in constant electromagnetic backgrounds, or suitable
generalizations thereof, this symmetry may be relevant to
noncommutative field theories (as already noted in
\cite{Doplicher:1994tu}). By the same token, its relevance to the
quantum Hall effect cannot be ruled out.}

\corr{Setting these speculations aside, in this paper we investigate a
related problem: the construction and properties of theories with
local Maxwell invariance. In Sec. II, following
\cite{Durka:2011nf}, we review the formulation of such a theory and
derive its field equations. Sec. III briefly discusses possible
relations between gauged Maxwell theories and bi-metric theories of
gravity. Finally, in Sec. IV, we examine a simple class of
cosmological solutions.}

\section{Gauged Maxwell theory}



We begin this section presenting some definitions and then recalling

some of the results presented in \cite{Durka:2011nf}.



The starting point of our construction is a generalization of

Maxwell symmetry, in which the Poincar\'e subalgebra of the original

Maxwell algebra \cite{Bacry:1970ye}, \cite{Schrader:1972zd} is

replaced with the anti de Sitter counterpart. In other words we

consider the algebra of symmetries of (complex scalar, Dirac) fields

on anti de Sitter space and constant electromagnetic

background\footnote{A detailed discussion will be presented in a

forthcoming paper.}. The algebra of such symmetries, which turns out

to be  isomorphic to $\sf{so}(3,1)\oplus\sf{so}(3,2)$ was discussed

previously in \cite{Soroka:2006aj}, \cite{Bonanos:2010fw},

\cite{Lukierski:2010dy} and has the following form

\cite{Durka:2011nf}

\begin{align}

&[\mathcal{P}_{a},\mathcal{P}_{b}]=i(\mathcal{M}_{ab}-\mathcal{Z}_{ab})\,

,\nonumber\\

&[\mathcal{M}_{ab},\mathcal{M}_{cd}]=-i(\eta_{ac}\mathcal{M}_{bd}+\ldots)\, ,\nonumber\\

&[\mathcal{M}_{ab},\mathcal{Z}_{cd}]=-i(\eta_{ac}\mathcal{Z}_{bd}+\ldots)\,,\nonumber\\

&[\mathcal{Z}_{ab},\mathcal{Z}_{cd}]=-i(\eta_{ac}\mathcal{Z}_{bd}+\ldots)\, ,\nonumber\\

&[\mathcal{M}_{ab},\mathcal{P}_c]=-i(\eta_{ac}\mathcal{P}_{b}-\eta_{bc}\mathcal{P}_{a}),\quad

[\mathcal{Z}_{ab},\mathcal{P}_{c}]=0\,,\label{1}

\end{align}

where $\ldots$ denotes three more terms obtained by

antisymmetrization in the pairs of indices $(ab)$ and $(cd)$. The

original Maxwell algebra \cite{Soroka:2006aj},

\cite{Bonanos:2010fw}, \cite{Lukierski:2010dy} can be easily

obtained from (\ref{1}) by contraction, which makes $\mathcal{M}$ on

the right hand side of the first commutator and the whole right hand

side of the fourth one vanish.



This algebra can be gauged by defining the gauge field (connection)

\begin{equation}\label{2}

\mathbb{A}_{\mu}=\frac{1}{2}\omega_\mu{}^{ab}

\mathcal{M}_{ab}+\frac{1}{\ell}e_\mu{}^{a}

\mathcal{P}_{a}+\frac{1}{2}h_\mu{}^{ab} \mathcal{Z}_{ab}\,,

\end{equation}

and its curvature

\begin{equation}\label{3}

\mathbb{F}_{\mu\nu}=\partial_\mu\mathbb{A}_{\nu}-\partial_\nu\mathbb{A}_{\mu}-i[\mathbb{A}_{\mu},\mathbb{A}_{\nu}]\,,

\end{equation}

which can be decomposed into Lorentz, translational, and Maxwell

parts

\begin{equation}\label{4}

\mathbb{F}_{\mu\nu}=\frac12\, F_{\mu\nu}^{ab}\, \mathcal{M}_{ab}

+\frac{1}{\ell}T_{\mu\nu}^{a}\,\mathcal{P}_{a}+ \frac12\,

G_{\mu\nu}^{ab}\, \mathcal{Z}_{ab}\,,

\end{equation}

where

\begin{eqnarray}

F_{\mu\nu}^{ab}&=&R^{ab}_{\mu\nu}+\frac{1}{\ell^2}( e^a_\mu e^b_\nu- e^a_\nu e^b_\mu) \,,\label{5a}\\

T_{\mu\nu}^{a}&=&D^\omega_\mu e^a_\nu -D^\omega_\nu e^a_\mu\,,\label{5b}\\

G_{\mu\nu}^{ab}&=&D^\omega_\mu h^{ab}_\nu -D^\omega_\nu

h^{ab}_\mu\nonumber-\frac{1}{\ell^2}( e^a_\mu e^b_\nu- e^a_\nu

e^b_\mu)\\ &+&(h^{ac}_\mu h^{\quad b}_{\nu\,c}-h^{ac}_\nu h^{\quad

b}_{\mu\,c})\, .\label{5c}

\end{eqnarray}

\corr{Here, $D^\omega_\mu(\ast)\equiv
\partial_\mu(\ast)-i[\omega_\mu,(\ast)]$ denotes the covariant
derivative associated with the Lorentz connection $\omega$.}



\corr{The gauge transformations of the gauge fields
$h_\mu{}^{ab}$, $\omega_\mu{}^{ab}$, and $e_\mu{}^a$ are}
\begin{eqnarray}

        \delta_\Theta h^{ab}_{\mu}&=&D^\omega_\mu \tau^{ab}+

h^{ac}_\mu(\lambda+\tau)_{c}^{\;\,b}+h^{bc}_\mu(\lambda+\tau)^{a}_{\;\;c}\label{6}\quad\\

      \delta_\Theta \omega^{ab}_{\mu}&=& D^\omega_\mu \lambda^{ab}

      \label{7}\\

   \delta_\Theta e^{a}_{\mu}&=&-\lambda^a_{\;b}

    \,e^b_\mu\,,\label{8}

\end{eqnarray}

\corr{where $\lambda^{ab}$ and $\tau^{ab}$ are the parameters of the
local Lorentz and Maxwell symmetries, respectively. The curvature
components transform homogeneously:}
\begin{eqnarray}

        \delta_\Theta G^{ab}_{\mu\nu}&=&-[\tau,F_{\mu\nu}]^{ab}-[(\lambda+\tau),G_{\mu\nu}]^{ab}\label{9}\\

      \delta_\Theta F^{ab}_{\mu\nu}&=&-[\lambda,F_{\mu\nu}]^{ab}  \label{10}\\

    \delta_\Theta

    T^{a}_{\mu\nu}&=&-\lambda^a_{\;b}T^b_{\mu\nu}\,.\label{11}

\end{eqnarray}

\corr{We have omitted local translations from the transformations above.
As usual, local translations are ultimately traded for general
coordinate invariance and are not symmetries of the action.}

\corr{It follows from Eqs. (\ref{8})--(\ref{11}) that the only
geometrical action that can be constructed from the gauge fields and
curvatures as a Lorentz- and Maxwell-invariant four-form (see
\cite{Durka:2011nf} and references therein) is}
\begin{equation}

S_{E} =\frac1{64\pi G}\,

\int\epsilon_{abcd}\left(R_{\mu\nu}{}^{ab}e_{\rho}^{c}e_{\sigma}^{d}

-\frac{\Lambda}{3}e_{\mu}^{a}e_{\nu}^{b}e_{\rho}^{ c}e_{\sigma}^{

d}\right)\epsilon^{\mu\nu\rho\sigma}\label{12}

\end{equation}

\corr{This is the standard Einstein action with a cosmological-constant
term. Other invariant terms quadratic in the curvatures are also
possible; however, these terms are topological and do not affect the
equations of motion.}

\corr{It is useful to compare this construction with the one presented
in \cite{deAzcarraga:2010sw}. There, the authors construct curvatures
similar to those in Eqs. (\ref{5a})--(\ref{5c}), but build the action
from all possible local-Lorentz-invariant geometric combinations.
The resulting theory therefore breaks Maxwell symmetry explicitly.
Consequently, the geometrical action in \cite{deAzcarraga:2010sw} is
more general than the one derived in \cite{Durka:2011nf} and reviewed
here.}

\corr{The action (\ref{12}) contains no curvature term for the Maxwell
field strength $G$: as noted above, the geometrical terms involving
$G$ are topological. It is therefore natural to ask whether
nontrivial terms involving the Maxwell field strength can be added.
An inspection of Eqs. (\ref{9}) and (\ref{10}) shows that two such
terms are possible. The first is a Yang--Mills action for the
$\sf{SO(3,1)}$ gauge group, already identified in
\cite{Soroka:2011tc} and \cite{Durka:2011nf}:}
\begin{equation}\label{13}

 S_{1}= -  \frac1{4\varrho_1}\,\int d^4x\, e\, \big( F_{\mu\nu}{}^{ab} + G_{\mu\nu}{}^{ab}\big) \big( F^{\mu\nu}{}_{ab} +

    G^{\mu\nu}{}_{ab}\big)\, .

\end{equation}

\corr{Because $\epsilon^{abcd}$ is an invariant tensor of
$\sf{SO(3,1)}$, we can construct a second manifestly gauge-invariant
action:}
\begin{equation}\label{13a}

 S_{2}= -  \frac1{4\varrho_2}\,\int d^4x\, e\, \big( F_{\mu\nu}{}^{ab} + G_{\mu\nu}{}^{ab}\big) \big( F^{\mu\nu}{}^{cd} +

    G^{\mu\nu}{}^{cd}\big) \epsilon_{abcd}\, .

\end{equation}

\corr{The parameters $\varrho_1$ and $\varrho_2$ are dimensionless
coupling constants. Using the explicit expressions for the curvatures
$F$ and $G$ in Eqs. (\ref{5a}) and (\ref{5c}), we can simplify the
actions further. In particular,}
$$
F_{\mu\nu}{}^{ab}(\omega, e) + G_{\mu\nu}{}^{ab}(h,\omega,e) =
H_{\mu\nu}{}^{ab}(\varpi)
$$
\corr{where $H_{\mu\nu}{}^{ab}(\varpi)$ is the curvature tensor of the
connection $\varpi_{\mu}{}^{ab}\equiv
\omega_{\mu}{}^{ab}+h_{\mu}{}^{ab}$. We take $\varpi$ as the new
fundamental variable. The total action is the sum of the Einstein
action (\ref{12}) and the actions for the field $\varpi$:}
\begin{align}

S &= S_E + S_1 +S_2 \nonumber\\

&=\frac{1}{64\pi

G}\int\epsilon_{abcd}\left(R_{\mu\nu}{}^{ab}e_{\rho}^{c}e_{\sigma}^{d}

-\frac{\Lambda}{3}e_{\mu}^{a}e_{\nu}^{b}e_{\rho}^{ c}e_{\sigma}^{

d}\right)\epsilon^{\mu\nu\rho\sigma}\nonumber\\

&-\frac14\,\int  e\, \left(\frac1{\varrho_1}\, H_{\mu\nu}{}^{ab}

H^{\mu\nu}{}_{ab}+ \frac1{\varrho_2}\,H_{\mu\nu}{}^{ab}

H^{\mu\nu}{}^{cd}\,\epsilon_{abcd}\right)\, .\label{14}
\end{align}
\add{The decomposition above is useful because it isolates the dynamical
connection $\varpi$ while retaining the gravitational variables in the
Einstein--Hilbert sector. It also makes clear that the Maxwell-sector
couplings $\varrho_1$ and $\varrho_2$ control distinct quadratic
invariants; their effects should therefore be assessed together rather
than by examining either term in isolation.}




\corr{Varying this action with respect to the tetrad gives the Einstein
equations with a cosmological-constant term:}
\begin{equation}\label{15a}

    R_{\mu\nu}-\frac12\, g_{\mu\nu}\, R +\Lambda\, g_{\mu\nu}= 8\pi G T_{\mu\nu}

\end{equation}

\corr{where $T_{\mu\nu}$ is the energy--momentum tensor of the field
$H$:}
\begin{align}

    T_{\mu\nu} &= \frac{1}{\varrho_1}\left(H_{\mu\lambda}{}^{ab} H_\nu{}^\lambda{}_{ab} -

    \frac14\, g_{\mu\nu}\,H_{\lambda\sigma}{}^{ab}

    H^{\lambda\sigma}{}_{ab}\right) \nonumber\\

    &+\frac{1}{\varrho_2}\left(H_{\mu\lambda}{}^{ab} H_\nu{}^\lambda{}^{cd} -

    \frac14\, g_{\mu\nu}\,H_{\lambda\sigma}{}^{ab}

    H^{\lambda\sigma}{}^{cd}\right)\epsilon_{abcd}\label{15}

\end{align}





\corr{The field equation for $\varpi$ differs slightly from the
standard Yang--Mills equation and takes the form}
\begin{equation}\label{16}

    \nabla^\mu \left(\frac1{\varrho_1}\, H_{\mu\nu}{}^{ab} +\frac1{\varrho_2}\,\epsilon^{abcd}

    H_{\mu\nu}{}_{cd}\right)=0\,,

\end{equation}

\corr{Here, $\nabla$ combines the Levi-Civita covariant derivative of
the metric constructed from the tetrad $e^a_\mu$ with the action of
the connection $\varpi$ on the algebra indices. The field equations
following from action (\ref{14}) imply that spacetime torsion
vanishes.}

\corr{Gauge theories with noncompact gauge groups are known to face
unitarity problems at the quantum level \cite{Hsu:1980zt}. Moreover,
even classically, coupling such theories to gravity is problematic:
the energy density is not positive definite, as can be seen from the
$T_{00}$ component of the energy--momentum tensor in Eq. (\ref{15}).
We examine this issue in greater detail in Sec. IV, where we study
simple cosmological solutions of the theory defined by action
(\ref{14}).}

\section{Bi-metric interpretation}



\corr{This section examines the relation between gauged Maxwell
theory and bi-metric theories. The gauge transformations of the
connection $\varpi$, introduced in the preceding section, have
formally the same structure as those of the gravitational connection
$\omega$. This observation allows us to introduce a second
``tetrad'' one-form, $f^a$, together with its associated
``torsion'' two-form, $\tau^a$:}
\begin{equation}\label{17}

    df^a + \varpi^a{}_b\wedge f^b=\tau^a\,.

\end{equation}

\corr{Once the new tetrad field $f^a$ has been defined, we can
construct another gauge-invariant action: the Einstein action
expressed in terms of $f$ and $\varpi$. In particular,}
\begin{equation}

S_{E'} =\frac1{64\pi G'}\,

\int\epsilon_{abcd}\left(H_{\mu\nu}{}^{ab}f_{\rho}^{c}f_{\sigma}^{d}

-\frac{\Lambda'}{3}f_{\mu}^{a}f_{\nu}^{b}f_{\rho}^{ c}f_{\sigma}^{

d}\right)\epsilon^{\mu\nu\rho\sigma}\label{18}

\end{equation}

\corr{this action is manifestly gauge invariant, just as the
Einstein action (\ref{12}) is. The sum $S_E+S_{E'}$ closely
resembles the $f$--$g$ theory of Isham, Salam, and Strathdee
\cite{Isham:1971gm} (see, e.g., \cite{Banados:2008fi} for a
more recent discussion) in the first-order formalism. The remaining
ingredient is a contact term coupling the tetrads $e$ and $f$.
Because these tetrads transform under different gauge groups, a
direct coupling between them is not available; instead, one can
couple the metrics constructed from them,
$g_{\mu\nu}=e_\mu{}^a e_\nu{}^b\eta_{ab}$ and
$q_{\mu\nu}=f_\mu{}^a f_\nu{}^b\eta_{ab}$.}

\corr{In the usual $f$--$g$ framework, such interaction terms reduce the
$(\mathrm{diff})^2$ invariance of $S_E+S_{E'}$ to its diagonal
subgroup. Here, however, the theory has only one built-in
diffeomorphism invariance, inherited from the gauged translations.
Adding a contact term therefore does not break an additional
symmetry. The most general such term can be written as}
\begin{equation}\label{19}

    S_{cont} =

   {M^2}(-g)^u(-q)^{(1/2-u)}\, V(g^{-1}q)

\end{equation}

\corr{where $(g^{-1}q)^\mu{}_\nu=g^{\mu\rho}q_{\rho\nu}$ and
$V$ is an arbitrary scalar potential constructed from this tensor.}
\add{This expression is a symmetry-based parametrization, not a claim
that every choice of $u$ and $V$ yields a dynamically healthy theory.
A viable model would require an independent study of the constraint
structure, propagating degrees of freedom, and stability of perturbations
about the backgrounds of interest.}




\corr{Bi-metric theories have been studied from many perspectives;
see, for example, \cite{Banados:2008fi}, \cite{Damour:2002wu},
\cite{Damour:2002ws}, \cite{ArkaniHamed:2002sp},
\cite{Blas:2005yk}, and \cite{Speziale:2010cf}, as well as the
comprehensive review \cite{Clifton:2011jh}. Determining whether the
local Maxwell symmetry has a natural realization in any of these
settings---including brane-world scenarios, Kaluza--Klein models,
or noncommutative geometry---is an interesting direction for future
work.}
\section{Example: a simple cosmological model}



\corr{We now consider a simple cosmological solution of the model
defined by action (\ref{14}). This example illustrates the
difficulty associated with the sign-indefinite energy density
discussed in Sec. II.}

\corr{We adopt the spatially flat Friedmann--Robertson--Walker (FRW)
metric:}
\begin{equation}\label{20}

    ds^2 =-N(t)^2dt^2+a(t)^2 (dx^2+dy^2+dz^2)

\end{equation}
\corr{We analyze the dynamics using the variational principle rather
than starting directly from the equations of motion. Substituting
the metric (\ref{20}) into the Einstein action (\ref{12}) and
omitting the (infinite) spatial-volume factor gives}
\begin{equation}

S_E=\frac{1}{8\pi G}\int dt \left(-\frac{3}{N}a\dot a^2-\Lambda a^3

N\right)\,.\label{21}

\end{equation}



\corr{We next consider the gauge-field contribution to the action.
The first step is to determine the most general field configuration
compatible with the background symmetries of the FRW spacetime
(\ref{20}). The relevant invariant tensors are those of the
Euclidean group of spatial translations and rotations. Accordingly,
we take the only nonvanishing components of the connection
$\varpi$ to be}
\begin{equation}\label{22}

    \varpi_i{}^{0\alpha} = f(t)\, \delta_i^\alpha\,,

    \quad \varpi_i{}^{\alpha\beta} = g(t)\,

    \epsilon_i{}^{\alpha\beta}\,,\quad \alpha,\beta=1,2,3\,.

\end{equation}

\corr{We derive this ansatz explicitly in the Appendix.}



\corr{Substituting this ansatz and the metric (\ref{20}) into the
action $S_1$ in (\ref{13}) gives}
\begin{equation}
    S_1 =\frac1{\varrho_1}\int dt \left( \frac{3a}{N}\Big(\dot{g}^2 -\dot{f}^2\Big)
+ \frac{3N}{a} g^2\Big(4f^2-g^2\Big) \right)\label{23}
\end{equation}
\corr{Similarly, the action $S_2$ in (\ref{14}) becomes}
\begin{equation}
    S_2 =\frac1{\varrho_2}\int dt \left( \frac{12a}{N}\, \dot f\,\dot g - \frac{24N}{a}\, f\, g^3\right)\label{24}
\end{equation}
\corr{The action (\ref{23}) contains a kinetic term with the
opposite sign from the conventional one, and its potential is not
bounded from below. These features illustrate the difficulties
associated with a noncompact gauge group.}

\corr{The Einstein equations follow from the total action
$S_E+S_1+S_2$ by varying with respect to the scale factor $a(t)$
and the lapse function $N(t)$. Since the lapse acts as a Lagrange
multiplier, we may set $N(t)=1$ after variation. The resulting
equations are}
\begin{align}

\left(\frac{\dot{a}}{a} \right)^2-\frac{\Lambda}{3}

 &=\frac{8\pi G}{3\varrho_1}\left[\frac{3}{a^2} \left(\dot g^2 - \dot

 f^2\right)

 -\frac{3}{a^4}\,g^2\left(4f^2-g^2\right)\right]\nonumber\\

 &+\frac{8\pi G}{3\varrho_2}\left[\frac{12}{a^2}\, \dot f \dot g +\frac{24}{a^4}\,f\, g^3\right]\,.\label{25}

\end{align}

and

\begin{align}

\frac{2\ddot a}{a} &+ \left(\frac{\dot a}{a}\right)^2 - \Lambda

 =-\frac{8\pi G}{3\varrho_1}\left[\frac{3}{a^2} \left(\dot g^2 - \dot

 f^2\right)\right.\nonumber\\

 &-\left.\frac{3}{a^4}\,g^2\left(4f^2-g^2\right)\right]

 -\frac{8\pi G}{3\varrho_2}\left[\frac{12}{a^2}\, \dot f \dot g +\frac{24}{a^4}\,f\, g^3\right]\label{26}

\end{align}

\corr{To identify the effective energy density and pressure, we
compare these equations with the standard Friedmann equations:}
$$

\left(\frac{\dot{a}}{a} \right)^2-\frac\Lambda3=\frac{8\pi

G}{3}\varepsilon\, , \quad \ddot a -\frac{\Lambda a}3=-\frac{4\pi

G}{3}(\varepsilon+3p)a

$$
\corr{The effective energy density is therefore}


\del{we see that in our case the energy density is}

\begin{align}

    \varepsilon&= \frac{1}{\varrho_1}\left[\frac{3}{a^2} \left(\dot g^2 - \dot

 f^2\right)

 -\frac{3}{a^4}\,g^2\left(4f^2-g^2\right)\right]\nonumber\\

 &+\frac{1}{\varrho_2}\left[\frac{12}{a^2}\, \dot f \dot g +\frac{24}{a^4}\,f\, g^3\right]\label{27}

\end{align}


\corr{The resulting energy density is not positive definite.}
\add{This sign-indefinite expression characterizes the reduced matter
sector before the gravitational constraint is imposed. Its interpretation
should therefore be distinguished from the positivity properties of
the total constrained cosmological solutions.}

\corr{The pressure is}

\begin{equation}\label{28}

    p=\frac\varepsilon3

\end{equation}

\corr{This equation of state is expected for a massless Yang--Mills
field. Thus, Eqs. (\ref{25}) and (\ref{26}) describe a universe
containing radiation and a cosmological constant. The corresponding
scale-factor solutions are}
\begin{align*}

 &a(t)= C \sqrt{2 (t-t_0)}\,, &\Lambda=0\\

 &a(t)= C \sqrt{\cos \left(\sqrt{\frac{|\Lambda|}{3} } \left(2 (t-t_0)\right)\right)}\,, &\Lambda<0\\

&a(t)= C \sqrt{\cosh \left(\sqrt{\frac{|\Lambda|}{3} } \left(2

(t-t_0)\right)\right)}\,, &\Lambda>0\,.

\end{align*}

\corr{The first Friedmann equation (\ref{25}) constrains the allowed
energy density: for negative or vanishing $\Lambda$, it requires a
nonnegative density, while for positive $\Lambda$ it imposes a lower
bound. The gravitational constraint therefore appears to exclude
some configurations of $\varpi$ with sufficiently large negative
energy density. Whether this restriction is a consequence of the
cosmological ansatz (\ref{20}) or a more general feature of the
theory remains an open question.}

For completeness we write the $g$ and $f$ equations of motion. They

read

\begin{equation}

\frac1{\varrho_1}\left[ -(a\dot g)\dot{} + \frac{4}a\, gf^2 -

\frac{2}a\, g^3\right] -\frac1{\varrho_2}\left[ (2a\dot f)\dot{} +

\frac{12}a\, fg^2 \right]=0\label{29}

\end{equation}

and

\begin{equation}

\frac1{\varrho_1}\left[ (a\dot f)\dot{} + \frac{4}a\, g^2f \right]

-\frac1{\varrho_2}\left[ (2a\dot g)\dot{} + \frac{4}a\, g^3

\right]=0\label{30}

\end{equation}

It is worth noticing that these equations are related to the

continuity equation

\begin{equation}\label{31}

\dot \varepsilon +3H(\varepsilon+p)= \dot \varepsilon

+4H\varepsilon=0\,,

\end{equation}

where we used (\ref{28}). Indeed substituting (\ref{27}) to

(\ref{31}) one finds

\begin{equation}\label{32}

    \dot f\, \mbox{\framebox{(\ref{29})}} +\dot g\, \mbox{\framebox{(\ref{30})}}=0\,.

\end{equation}







\corr{For the simplest case, in which $f$ and $g$ are
time-independent constants, we exclude the trivial solution $g=0$.
Equation (\ref{30}) then gives}
$$f=\frac{\varrho_1}{\varrho_2}\,g\,.$$
\corr{Substitution into Eq. (\ref{29}) shows that a nontrivial
solution would require}
$$\left(\frac{\varrho_1}{\varrho_2}\right)^2=-\frac1{4}\,,$$
\corr{which has no solution for real coupling constants. A similar
conclusion follows under the assumption
$\dot g,\dot f\sim a^{-1}$. These observations suggest that the
dynamics of $f$ and $g$ is nontrivial and merits further study.}

\section{Concluding remarks}

\add{We have reviewed a construction in which local Maxwell symmetry is
incorporated into a gravitational gauge theory and examined several
consequences of the resulting action. The formulation admits a
Yang--Mills-like description in terms of a composite connection, while
also suggesting a relation to bi-metric models through the introduction
of a second tetrad. A homogeneous and isotropic reduction provides a
tractable setting in which to inspect the cosmological dynamics and the
sign-indefinite energy contributions associated with the noncompact
gauge group.}

\add{The results presented here should be regarded as a preliminary
classical analysis. In particular, the cosmological ansatz does not
establish the stability of general inhomogeneous solutions, and the
bi-metric interpretation does not settle the number or health of the
propagating degrees of freedom. Further work should address perturbative
stability, the constraint algebra, and the extent to which admissible
parameter choices can avoid the pathologies associated with
noncompact Yang--Mills sectors. These questions are necessary before
the framework can be assessed as a candidate for a physically viable
extension of gravity.}

\appendix*

\section{Cosmological gauge field}



\corr{In this appendix, we derive the most general form of the gauge
field $h_\mu^{ab}$ on the spatially flat FRW background. The
infinitesimal symmetries of this background are spatial translations
and rotations, which together generate the three-dimensional
Euclidean group. Since the spatial metric on each constant-time
hypersurface is proportional to $\delta_{ij}$, we do not distinguish
carefully between upper and lower spatial indices in the following.}
\begin{equation}\label{a1}

    [R_i, R_j] = \epsilon_{ijk}\, R_k, \, [R_i, T_j] = \epsilon_{ijk}\, T_k, \, [T_i, T_j]

    =0\, .

\end{equation}

\corr{The Killing vector fields of the spatially flat FRW spacetime
provide a representation of this algebra:}
\begin{align}

T_i &= \frac{\partial}{\partial x^i}\nonumber \\ R_i

&=\epsilon_{ijk}\, x^j \,\frac{\partial}{\partial x^k}\label{a2}

\end{align}



\corr{Invariance of the gauge field under spacetime symmetries means
that its Lie derivative along each symmetry-generating Killing
vector vanishes up to a gauge transformation (for general
treatments, see \cite{Forgacs:1979zs} and
\cite{Jackiw:1979ub}):}
\begin{align}

    {\cal L}_{T_i}\,h_\mu^{ab} &= \partial_\mu \lambda^{(T)}_i{}^{ab}

    + [h_\mu,\lambda^{(T)}_i]^{ab}\,,\label{a3}\\  {\cal L}_{R_i}\,h_\mu^{ab} &= \partial_\mu \lambda^{(R)}_i{}^{ab}

    + [h_\mu,\lambda^{(R)}_i]^{ab}\,.\label{a4}

\end{align}

\corr{These equations obey integrability conditions that follow from
the identity}
$$

[{\cal L}_X, {\cal L}_Y] = {\cal L}_{[X,Y]}\, ,

$$

\corr{which yields the following condition on the gauge parameters}

$\lambda$

\begin{equation}\label{a5}

{\cal L}_X\lambda^{(Y)}-{\cal L}_Y\lambda^{(X)}

-[\lambda^{(X)},\lambda^{(Y)}] =\lambda^{([X,Y])}\,.

\end{equation}

\corr{For two translations, Eq. (\ref{a5}) becomes}
\begin{equation}\label{a6}

    \partial_i\lambda^{(T)}_j{}^{ab}-

    \partial_j\lambda^{(T)}_i{}^{ab} -

    [\lambda^{(T)}_i,\lambda^{(T)}_j]^{ab}=0\, .

\end{equation}

\corr{This equation states that the field strength of the
translation-associated gauge parameters vanishes. Since Eq.
(\ref{a3}) is gauge covariant, we may choose the gauge in which
$\lambda^{(T)}_j{}^{ab}=0$.}

\corr{Applying Eq. (\ref{a5}) to a rotation and a translation then
gives}
$$

\partial_i\lambda^{(R)}_j{}^{ab}=0

$$

\corr{so $\lambda^{(R)}_j{}^{ab}$ is constant.} Then the

equation (\ref{a5}) applied to rotations takes the form

\begin{equation}\label{a7}

    [\lambda^{(R)}_i,\lambda^{(R)}_j]^{ab}=\epsilon_{ijk}\,

    \lambda^{(R)}_k{}^{ab}\,,

\end{equation}

which means that the matrices $\lambda^{(R)}_i{}^{ab}$ form a

representation of the rotation group $\sf{SO(3)}$.



\corr{With this result, we can return} to the equations (\ref{a3}),

(\ref{a4}) which now take a much simpler form

\begin{align}

    {\cal L}_{T_i}\,h_\mu^{ab} &= 0\label{a8}\\  {\cal L}_{R_i}\,h_\mu^{ab} &=  [h_\mu,\lambda^{(R)}_i]^{ab}\label{a9}

\end{align}

\corr{Equation (\ref{a8}) shows that} the gauge field depends only on time

$h_\mu^{ab}(x,t) = h_\mu^{ab}(t)$. Then from (\ref{a9}) it follows

that $h_0^{ab}$ commutes with all the matrices $\lambda^{(R)}_i$ and

therefore must be zero, and

\begin{equation}\label{a10}

\epsilon_{ijk}\, h_j^{ab}=[h_k,\lambda^{(R)}_i]^{ab}

\end{equation}

Returning to eq.\ (\ref{a7}) we see that the matrices

$\lambda^{(R)}_i{}^{ab}$ could be represented as matrices of

generators of rotations in the Lorentz group $\sf{SO(3,1)}$ so that

explicitly

\begin{equation}\label{a11}

\lambda^{(R)}_i{}^{\alpha\beta} = \epsilon_i{}^{\alpha\beta}\,,

\quad\lambda^{(R)}_i{}^{\alpha0} =0\,, \quad \alpha,\beta=1,2,3

\end{equation}

Substituting this in (\ref{a10}) one can easily check that there are

two solutions

\begin{equation}\label{a12}

h_j^{0\beta} =f(t) \, \delta_j^\beta\,, \quad h_j^{\alpha\beta} =

g(t)\, \epsilon_i{}^{\alpha\beta}

\end{equation}

\corr{This is the ansatz adopted in the main text.}



\del{[Removed: the original acknowledgments refer to correspondence and grants associated with the source authors, not with the fictional authors used in this adaptation.]}









































\begin{thebibliography}{99.}



%\cite{Bacry:1970ye}

\bibitem{Bacry:1970ye}

  H.~Bacry, P.~Combe, J.~L.~Richard,

  ``Group-theoretical analysis of elementary particles in an external electromagnetic field. 1.

  the relativistic particle in a constant and uniform field,''

  Nuovo Cim.\  {\bf A67 } (1970)  267-299.



%\cite{Schrader:1972zd}

\bibitem{Schrader:1972zd}

  R.~Schrader,

  ``The maxwell group and the quantum theory of particles in classical homogeneous electromagnetic fields,''

  Fortsch.\ Phys.\  {\bf 20 } (1972)  701-734.



  %\cite{Galindo:1967}

\bibitem{Galindo:1967}

  A.~Galindo,

  ``Lie algebra extensions of the Poincar\'e Algebra,''

  J.\ Math.\ Phys.\  {\bf 8 } (1967)  768.



%\cite{Soroka:2011tc}

\bibitem{Soroka:2011tc}

  D.~V.~Soroka and V.~A.~Soroka,

  ``Gauge semi-simple extension of the Poincar\'e group,''

  arXiv:1101.1591 [hep-th].

  %%CITATION = ARXIV:1101.1591;%%



%\cite{Gomis:2009dm}

\bibitem{Gomis:2009dm}

  J.~Gomis, K.~Kamimura and J.~Lukierski,

  ``Deformations of Maxwell algebra and their Dynamical Realizations,''

  JHEP {\bf 0908}, 039 (2009)

  [arXiv:0906.4464 [hep-th]].

  %%CITATION = JHEPA,0908,039;%%



  %\cite{Douglas:2001ba}

\bibitem{Douglas:2001ba}

  M.~R.~Douglas, N.~A.~Nekrasov,

  ``Noncommutative field theory,''

  Rev.\ Mod.\ Phys.\  {\bf 73 } (2001)  977-1029.

  [hep-th/0106048].



  %\cite{Doplicher:1994tu}

\bibitem{Doplicher:1994tu}

  S.~Doplicher, K.~Fredenhagen, J.~E.~Roberts,

  ``The Quantum structure of space-time at the Planck scale and quantum fields,''

  Commun.\ Math.\ Phys.\  {\bf 172 } (1995)  187-220.

  [hep-th/0303037].



%\cite{Bonanos:2009wy}

\bibitem{Bonanos:2009wy}

  S.~Bonanos, J.~Gomis, K.~Kamimura, J.~Lukierski,

 ``Maxwell Superalgebra and Superparticle in Constant Gauge Backgrounds,''

  Phys.\ Rev.\ Lett.\  {\bf 104}, 090401 (2010).

  [arXiv:0911.5072 [hep-th]].



%\cite{deAzcarraga:2010sw}

\bibitem{deAzcarraga:2010sw}

  J.~A.~de Azcarraga, K.~Kamimura, J.~Lukierski,

  ``Generalized cosmological term from Maxwell symmetries,''

  Phys.\ Rev.\  {\bf D83}, 124036 (2011).

  [arXiv:1012.4402 [hep-th]].













  %\cite{Soroka:2006aj}

\bibitem{Soroka:2006aj}

  D.~V.~Soroka, V.~A.~Soroka,

  ``Semi-simple extension of the (super)Poincare algebra,''

  Adv.\ High Energy Phys.\  {\bf 2009 } (2009)  234147.

  [hep-th/0605251].



  %\cite{Bonanos:2010fw}

\bibitem{Bonanos:2010fw}

  S.~Bonanos, J.~Gomis, K.~Kamimura, J.~Lukierski,

  ``Deformations of Maxwell Superalgebras and Their Applications,''

  J.\ Math.\ Phys.\  {\bf 51 } (2010)  102301.

  [arXiv:1005.3714 [hep-th]].



  %\cite{Lukierski:2010dy}

\bibitem{Lukierski:2010dy}

  J.~Lukierski,

  ``Generalized Wigner-Inonu Contractions and Maxwell (Super)Algebras,''

  Proc.\ Steklov Inst.\ Math.\  {\bf 272} (2011) 1

  [arXiv:1007.3405 [hep-th]].

  %%CITATION = PSMMB,272,1;%%



  %\cite{Durka:2011nf}

\bibitem{Durka:2011nf}

  R.~Durka, J.~Kowalski-Glikman, M.~Szczachor,

  ``Gauged AdS-Maxwell algebra and gravity,''

[arXiv:1107.4728 [hep-th]].







  %\cite{Hsu:1980zt}

\bibitem{Hsu:1980zt}

  J.~P.~Hsu, M.~D.~Xin,

  ``Incompatibility Of Unitarity And Gauge Symmetry In The Sl(2,c) Yang-mills Field Theory,''

  Phys.\ Rev.\  {\bf D24 } (1981)  471.



  %\cite{Isham:1971gm}

\bibitem{Isham:1971gm}

  C.~J.~Isham, A.~Salam, J.~A.~Strathdee,

  ``F-dominance of gravity,''

  Phys.\ Rev.\  {\bf D3 } (1971)  867-873.



  %\cite{Banados:2008fi}

\bibitem{Banados:2008fi}

  M.~Banados, A.~Gomberoff, D.~C.~Rodrigues, C.~Skordis,

  ``A Note on bigravity and dark matter,''

  Phys.\ Rev.\  {\bf D79 } (2009)  063515.

  [arXiv:0811.1270 [gr-qc]].



  %\cite{Damour:2002wu}

\bibitem{Damour:2002wu}

  T.~Damour, I.~I.~Kogan, A.~Papazoglou,

  ``Nonlinear bigravity and cosmic acceleration,''

  Phys.\ Rev.\  {\bf D66 } (2002)  104025.

  [hep-th/0206044].



  %\cite{Damour:2002ws}

\bibitem{Damour:2002ws}

  T.~Damour, I.~I.~Kogan,

 ``Effective Lagrangians and universality classes of nonlinear bigravity,''

  Phys.\ Rev.\  {\bf D66 } (2002)  104024.

  [hep-th/0206042].





%\cite{ArkaniHamed:2002sp}

\bibitem{ArkaniHamed:2002sp}

  N.~Arkani-Hamed, H.~Georgi, M.~D.~Schwartz,

  ``Effective field theory for massive gravitons and gravity in theory space,''

  Annals Phys.\  {\bf 305 } (2003)  96-118.

  [hep-th/0210184].



  %\cite{Blas:2005yk}

\bibitem{Blas:2005yk}

  D.~Blas, C.~Deffayet, J.~Garriga,

 ``Global structure of bigravity solutions,''

  Class.\ Quant.\ Grav.\  {\bf 23 } (2006)  1697-1719.

  [hep-th/0508163].



%\cite{Speziale:2010cf}

\bibitem{Speziale:2010cf}

  S.~Speziale,

  ``Bi-metric theory of gravity from the non-chiral Plebanski action,''

  Phys.\ Rev.\  {\bf D82 } (2010)  064003.

  [arXiv:1003.4701 [hep-th]].



%\cite{Clifton:2011jh}

\bibitem{Clifton:2011jh}

  T.~Clifton, P.~G.~Ferreira, A.~Padilla, C.~Skordis,

 ``Modified Gravity and Cosmology,''



  [arXiv:1106.2476 [astro-ph.CO]].









































%\cite{Forgacs:1979zs}

\bibitem{Forgacs:1979zs}

  P.~Forgacs, N.~S.~Manton,

  ``Space-Time Symmetries in Gauge Theories,''

  Commun.\ Math.\ Phys.\  {\bf 72 } (1980)  15.



%\cite{Jackiw:1979ub}

\bibitem{Jackiw:1979ub}

  R.~Jackiw, N.~S.~Manton,

  ``Symmetries And Conservation Laws In Gauge Theories,''

  Annals Phys.\  {\bf 127 } (1980)  257.



\end{thebibliography}





\end{document}

